% sec-03: Document 03, SB 964, model transparency and clinical data protection.
% STAGE 3.  Two column widths retuned; no content change.
\docpart{SB 964 Data Protection}{SB 964 - California Oncology Model Transparency
and Clinical Data Protection Act of 2026}

\section{Findings and Declarations}

\begin{enumerate}
\item California has legislated three times on artificial intelligence in
  health care and once on training data transparency, and none of the four
  instruments reaches an interventional clinical trial
  \cite{ab3030_2024,ab489_2025,sb1120_2024,ab2013_2024}.

\item AB 2013 requires a developer to publish a documentation summary of the
  data used to train a generative system \cite{ab2013_2024}. It does not require
  a clinical trial site to say which model version produced a given completion,
  and in a trial that distinction decides whether a result can be reconstructed.

\item A prompt and its completion, once a clinician acts on them, are source
  records. Federal law already requires source records to be attributable,
  legible, contemporaneous, original, and accurate, and to be retained
  \cite{cfr11,cfr312}. This act states how that obligation applies to a model.

\item A national platform specification published in 2026 established a
  hash-chained audit trail with provenance services as a workable construction
  for exactly this record, across sites rather than within one
  \cite{natplatform2026}.

\item De-identification obligations under federal health privacy law and the
  research use limits under the Common Rule both apply, and a site that meets one
  does not thereby meet the other \cite{cfr45part164,cfr45part46}.

\item A model version that cannot be named cannot be audited. This act therefore
  makes the version register, not the model, the object of regulation.
\end{enumerate}

\section{Definitions}

Terms defined in document 01 §2 carry the same meaning here. In addition:

\begin{enumerate}
\item ``Model'' means a large language model or an ensemble of such models
  operated on premises at a qualified site.

\item ``Model version'' means a specific set of weights, system prompt,
  retrieval corpus, decoding parameters, and guardrail configuration, identified
  by a version string and a content hash.

\item ``Prompt'' means the complete input presented to a model version,
  including retrieved context.

\item ``Completion'' means the complete output returned by a model version.

\item ``Hash-chained audit trail'' means an append-only record in which each
  entry carries the cryptographic hash of the preceding entry, so that deletion
  or alteration of any entry is detectable.

\item ``Provenance record'' means the linkage from a completion to the model
  version, the prompt, the retrieval sources, the requesting user, and the human
  disposition.

\item ``Secondary use'' means any use of trial data other than the conduct,
  analysis, reporting, or regulatory submission of the trial for which it was
  collected.
\end{enumerate}

\section{Model Transparency Obligations}

\begin{mvtable}
\begin{tabularx}{\textwidth}{>{\raggedright\arraybackslash}p{4.4cm} >{\raggedright\arraybackslash}p{2.6cm} >{\raggedright\arraybackslash}X}
\headrow \thead{Obligation} & \thead{Frequency} & \thead{Record that discharges it}\\
\midrule
Training data provenance summary & On registration, and on any change & The developer's published documentation summary, filed with the department \\
Model version register entry & Before first clinical use of a version & The register row: version string, content hash, effective date, approver \\
Retrieval corpus manifest & Quarterly & The manifest with a content hash per source document \\
Guardrail configuration statement & On each version & The configuration file hash and a plain-language statement of what is blocked \\
Known limitation statement & On each version, and on discovery & The limitation log, which carries the date of discovery and the mitigation \\
\end{tabularx}
\end{mvtable}
\tabcapl{tab:sb964obl}{Five transparency obligations, each discharged by a
record rather than by a policy. The version register is the same artifact the
federal change control plan requires under document 07 §5, not a second one.}

A qualified site shall not place a model version into clinical use before its
register entry exists. The register is the object the department surveys and the
object a federal reviewer is offered, and maintaining one artifact for both is a
requirement of this section rather than a convenience.

\section{Clinical Data Protection}

\begin{enumerate}
\item Trial data shall be de-identified under the expert determination or safe
  harbor method before any use outside the conduct of the trial
  \cite{cfr45part164}.

\item Secondary use is prohibited without specific, separately documented
  consent. A general consent to research does not authorize secondary use of
  prompts or completions.

\item A prompt or completion that contains identifiers is protected health
  information and is handled as such. A site shall not rely on the fact that a
  completion was generated rather than recorded to treat it as anything else.

\item Where the federal privacy standard and the Common Rule differ, the more
  protective governs, and the site shall document which standard it applied and
  why \cite{cfr45part164,cfr45part46}.
\end{enumerate}

\section{Audit Trail and Retention}

\begin{enumerate}
\item The audit trail shall record, for each model interaction: the model
  version string and content hash; the prompt; the completion; the retrieval
  sources; the requesting user; the timestamp to the second; and the human
  disposition, which is one of accepted, modified, or rejected, with the name of
  the clinician who made it.

\item The audit trail shall be append-only and hash-chained. A site shall
  verify chain integrity monthly and shall retain the verification record.

\item Records shall be retained for six years from database lock. Federal law
  requires two years after the later of marketing application approval or
  discontinuation of shipment with notice to the agency \cite{cfr312}, and state
  facility record law requires less. Where the periods differ, the longest
  governs, and six years is the period this chapter sets.

\item Access logs to any premises zone classified as restricted or machine under
  document 09 §1 shall be retained for the same six years, so that a question
  about who was present can be answered for as long as a question about what the
  model produced.
\end{enumerate}

\begin{mvtable}
\begin{tabularx}{\textwidth}{>{\raggedright\arraybackslash}p{4.2cm} >{\raggedright\arraybackslash}p{2.5cm} >{\raggedright\arraybackslash}X}
\headrow \thead{Record class} & \thead{Retention} & \thead{Authority, and why this period}\\
\midrule
Prompts, completions, dispositions & 6 years from lock & Longest of 21 CFR 312.62, state facility law, and this chapter \\
Model version register & 6 years from lock & Must outlive every completion it explains \\
Chain integrity verifications & 6 years from lock & Evidence that the trail was not altered \\
Premises access logs & 6 years from lock & Document 09 §2; the presence question must survive as long as the output question \\
Consent records & 6 years from lock & 21 CFR 50 as adapted \cite{cfr50adapt} \\
\end{tabularx}
\end{mvtable}
\tabcapl{tab:sb964ret}{Retention by record class. The period is uniform at six
years by design: a package whose classes expire on different dates produces an
archive that cannot answer a question about a single day.}

\section{Enforcement and Breach Notification}

\begin{enumerate}
\item A site shall notify the department of a breach affecting trial data within
  72 hours of discovery, and affected participants within 15 days. Notification
  under existing California breach law is not displaced by this section.

\item Deficiencies under this chapter are classified as Class A, Class B, or
  Class C with the meanings and penalty ranges given in document 01 §5. A
  failure of chain integrity that cannot be explained is a Class A deficiency.

\item A site that places a model version into clinical use without a register
  entry commits a Class B deficiency for the first occurrence and a Class A
  deficiency for any subsequent occurrence.
\end{enumerate}
